\documentclass[10pt,a4paper]{amsart}
\usepackage[margin=19mm]{geometry}
\usepackage{amsmath,amssymb,mathtools}
\usepackage[scr=rsfs,cal=euler]{mathalfa}
\usepackage{xfrac}
\usepackage[unicode,hidelinks,bookmarks=false]{hyperref}
\newcommand{\ie}{i.e.\ }
\newcommand{\cf}{cf.\ }
\newcommand{\resp}{respectively\ }
\newcommand{\Subsection}{\subsection}
\providecommand{\nobreakdash}{\nobreak\hbox{-}}
\setlength{\emergencystretch}{2em}
\multlinegap0pt
\usepackage{amssymb}
\usepackage{enumerate}
\newtheorem{lemma}{Lemma}
\title{On the separation \L ojasiewicz exponents of real analytic sets in the real plane}
\author{PHI-D{U}NG HO{A}NG$^\dag$, HONG-DUC NGUYEN$^\ddag$}
\date{Source: arXiv:2603.15532v2 (CC BY 4.0)}
\begin{document}
\maketitle

\noindent\emph{Source and licence.} On the separation \L ojasiewicz exponents of real analytic sets in the real plane. Version arXiv:2603.15532v2 (CC BY 4.0), distributed by arXiv under the Creative Commons Attribution 4.0 International licence, http://creativecommons.org/licenses/by/4.0/, as recorded in the arXiv licence metadata for this exact version and shown on the abstract page https://arxiv.org/abs/2603.15532v2. Full licence notice: \texttt{licenses/arxiv-2603.15532-CC-BY-4.0.txt}.\par\smallskip

\noindent\textit{Editorial scope.} Counted unit: the article's lemma claim1 (source0.tex lines 306--312). The article's complete original proof of it is delivered with it (source0.tex lines 313--361, one \texttt{proof} environment of 49 lines). No mathematics is written, repaired or re-derived: the statement and the proof are the article's own lines, in the article's own notation, and no retained line is substituted or re-flowed.  Retained uncounted as the source-local context the proof needs: lines 187--196.  Nothing was omitted from any retained span, so the package carries no omission record.  No retained line contains a citation, so the wrapper carries no bibliography and no bibliography entry.  No retained line contains a figure environment, an image-inclusion command or drawing code, so no image is delivered and the package declares no asset dependency.  Editorial formatting only, in addition: an AMS \texttt{amsart} preamble that restates the article's own declarations and macros used by the retained lines, namely the environments \texttt{lemma} on separate counters, together with the standard packages those lines need.  The retained environments print in this document as Lemma 1 in order of appearance, whereas the article prints the same items with the numbering of its own full document.  The complete native source of the article as distributed in the arXiv e-print for this exact version is bundled unchanged as \texttt{sources/196485/source0.tex}.
\begin{lemma}\label{distant-sliding}
Suppose that $\phi$ is not a root of $f = 0$. Then
	\begin{itemize}
		\item[(i)] $\mathrm{ord}d(\phi, V_f)\leq \tan \theta_{E_1}.$		
		\item[(ii)] Let $\psi$ be a series of form
	$$\psi : x = \phi(y) + c y^{\rho} + o(\rho),\ c \in \mathbb{R}.$$
Suppose that $\rho>\mathrm{ord}d(\psi, V_f)$, or $\rho\geq \mathrm{ord}d(\psi, V_f)$ and $c$ is generic. Then
		$$\mathrm{ord}d(\psi, V_f)= \mathrm{ord}d(\phi, V_f).$$
	\end{itemize}
\end{lemma}

	\begin{lemma}\label{claim1} 
		Suppose that $\ell(\phi) > \mathscr{L}_+(V_f, V_g)$. Then 
    \begin{itemize}
        \item[(i)] $\tan\theta_{F_1} = \tan\theta_{G_1}.$
        \item[(ii)] The polynomial $\mathcal{E}_{F_1}\mathcal{E}_{G_1}$ has only one root, which is a real value.
    \end{itemize}
		\end{lemma}

	\begin{proof}
    %It is easy to see that $\phi$ is the form 	$$ (x = \phi(t), y = t)\ \text{or}\ (x = \phi(t), y = -t)\ \text{where}\ \phi(t) \in \mathbb{R}\{t^{1/N}\}. $$ 	Without loss of generality, assume that $\phi$ can be parametrized by $(x = \phi(t), y = t)$ and $\tan\theta_{F_1} > \tan\theta_{G_1}$. 
    By contradiction we assume that $\tan\theta_{F_1} \neq \tan\theta_{G_1}.$ Without loss of generality suppose that $\tan\theta_{F_1} > \tan\theta_{G_1}.$ Let $\phi_{1,\infty}$ and $\phi_{2,\infty}$ be final results of sliding $\phi$ along $f$ and $g$, respectively. Then we have 
	\begin{align}
		\phi_{1,\infty}(y) &= \phi(y) + \sum_{i \ge 1}a_iy^{\alpha_i}, a_1\neq 0\\
		\phi_{2,\infty}(y) &= \phi(y) + \sum_{i \ge 1}b_iy^{\beta_i}, b_1\neq 0,
	\end{align}
	where $\tan\theta_{F_1}=\alpha_1 <\alpha_2 < \ldots$ and $\tan\theta_{G_1} =\beta_1 < \beta_2 <
    \ldots $.
    
    Let $\phi_{1,2}$ be the approximation of $\phi_{1,\infty}$ and $\phi_{2,\infty}$. Then $\phi_{1,2}\in \mathcal{V}_a(fg)$ and
	$$ \phi_{1,2}(y) = \phi(y) + cy^{\tan\theta_{G_1}} + o(\tan\theta_{G_1}), $$ where $c \in \mathbb{R}$ is generic. It follows from Lemma \ref{distant-sliding} that
    	\begin{align}\label{order-with-Vg}
    		\mathrm{ord}d(\phi_{1,2}, V_g) = \mathrm{ord}d(\phi, V_g) \text{ and } \mathrm{ord}d(\phi_{1,2}, V_h) = \mathrm{ord}d(\phi, V_h).
    	\end{align}
    	 
We will show that $\ell(\phi_{1,2})=\ell(\phi)$ by considering the following two cases.

\begin{itemize}
	\item Case 1: 	$\mathrm{ord}d(\phi, V_f) \leq \tan\theta_{G_1}$. Then $\mathrm{ord}\ d(\phi_{1,2}, V_f) = \mathrm{ord}\ d(\phi, V_f)$ due to Lemma \ref{distant-sliding}. Hence $\ell(\phi_{1,2})=\ell(\phi)$.
    
	\item Case 2: $\mathrm{ord}d(\phi, V_f) > \tan\theta_{G_1}$. Take a real Newton-Puiseux root $\gamma$ of $f$ satisfying
	\begin{align*}
		\mathrm{ord}(\phi - \gamma) = \mathrm{ord}\ d(\phi, V_f).
	\end{align*}
    Since
    $$ \phi_{1,2}(y)- \gamma(y) =\phi(y)-\gamma(y) + cy^{\tan\theta_{G_1}} + o(\tan\theta_{G_1}), $$
     it follows that $\mathrm{ord}d(\phi_{1,2} - \gamma) = \tan\theta_{G_1}$. Therefore,
	\begin{align}\label{order-with-Vf-tan}
		\mathrm{ord}d(\phi_{1,2}, V_f) \ge \tan\theta_{G_1}\ge \mathrm{ord}d(\phi_{1,2}, V_g).
	\end{align}
	Hence,
	  		\begin{align*}
			\ell(\phi_{1,2}) &= \frac{\min\{\mathrm{ord}d(\phi_{1,2},V_f), \mathrm{ord}d(\phi_{1,2},V_g)\}}{\mathrm{ord}d(\phi_{1,2},V_h)} \\
			&= \frac{ \mathrm{ord}d(\phi,V_g) }{\mathrm{ord}d(\phi,V_h)}  =\ell(\phi).
		\end{align*} 
		
\end{itemize}
The equality $\ell(\phi_{1,2})=\ell(\phi)$ contradicts the assumption that $\ell(\phi) > \mathscr{L}_+(V_f, V_g)$.  This proves (i).

(ii) By (i) one has $\rho=\tan \theta_{F_1} = \tan\theta_{G_1}$. Then $\mathcal{E}_{F_1}\mathcal{E}_{G_1}$ is the polynomial associated with the highest Newton edge of $\mathbb{P}(fg, \phi)$ with the angle $\rho$. Suppose by contradiction that $\mathcal{E}_{F_1}\mathcal{E}_{G_1}$ has two distinct roots $c_1,c_2$. Let $\varphi_i=\phi+c_i y^{\rho }$ for each $i=1,2$, and let $\varphi_{i,\infty}$ be a final result of sliding $\varphi_i$ along $fg$. Let $\varphi_{1,2}$ be the approximation of $\varphi_{1,\infty}$ and $\varphi_{2,\infty}$. Then $\varphi_{1,2}\in \mathcal{V}_a(fg)$ and 
	$$ \varphi_{1,2}(y) = \phi(y) + cy^{\rho} + o(\rho), $$ where $c \in \mathbb{R}$ is a generic number. Note that 
    $$\tan \theta_{H_1}\leq \tan \theta_{F_1}=\tan \theta_{G_1}=\rho.$$
    It follows from Lemma \ref{distant-sliding} that
    $$\mathrm{ord}d(\varphi_{1,2},V_f)=\mathrm{ord}d(\phi,V_f),\ \mathrm{ord}d(\varphi_{1,2},V_g)=\mathrm{ord}d(\phi,V_g) $$
	and
    $$d(\varphi_{1,2},V_h)=\mathrm{ord}d(\phi,V_h).$$
    Hence $\ell(\phi_{1,2})=\ell(\phi)$, which is a contradiction. The lemma is proved.    
\end{proof}

\end{document}
